\scanpage{49}
Le théorème de Thom est alors qu'un $x \handwritten{\in} H^{\cdot}(X)$ est dans S(X) sss il est
\deleted{stable}\inserted{fixe} par G, i.e. satisfait $A^{\handwritten{+}}. x = 0$. (Pourquoi cette condition est-el\deleted{\unsure{l}}%
le déjà nécessaire ?) U\handwritten{n}e autre façon de le voir est de dire que le
foncteur $H^{\cdot}$ considéré comme allant de la catégorie $\underline{S}$ précédente dans
la catégorie des \deleted{vectoriels} représentations de G dans des vectoriels
sur $\underline{F}_2$ , est pleinement fidèle. \deleted{Comme G \unsure{n'|est|et}} Quand on prend l'enveloppe
pseudo-abélienne, on trouve une sous-catégorie pleine de la catégorie des
représentations de G, mais il est douteux qu'on trouve une catégorie abé-%
lienne, voire la catégorie de toutes les représentations de G de dimensio
finie sur $\underline{F}_2$ . (Quelle est la structure de G, qui est prolisse; d'après
\margin{\handwritten{oui, car son anneau affine est une algèbre de polynômes à une infinité d'indéterminées, \deleted{\illegible{1}} réunion d'algèbres de polynômes à \unsure{engendrement} fini \unsure{stables} par $\Delta$}}Quillen, il serait pro-unipotent ....). Ici G bien sûr joue le rôle d'un
groupe de Galois motivique, mais contrairement à la situation familière,
il semble qu'il soit fortement non réductif (et d'ailleurs le corps de
définition dans tout ceci est $\underline{F}_2$ , et non $\underline{Q}$ ! ).

Quillent donne une généralisation de la construction d'une caté-%
gorie universelle $\underline{B}$ au cas où on travaille avec des variétés pas né\unsure{c|e}essai%
rement propres (mais f. n'étant supposé défini que pour f propre), et
quand on introduit une condition d'orientation \inserted{stable} sur les \uline{morphismes}
envisagés , via un morphisme $H \handwritten{\to BO(\infty)}$, où H est un H-espace (commuta%
tif, associatif et unitaire ?): une orientation de $f:X \handwritten{\to} Y$ sera un
relèvement homotopique de $X \handwritten{\to BO}$ défini par \deleted{$f^{\cdot}(T$} $T_X - f^{\cdot}(T_Y)$ en $X \handwritten{\to} H$.
les éléments de H(X,Y) seront définis par des diagrammes de morphismes
orientés $Z \handwritten{\to} X$, $Z \handwritten{\leftarrow} Y$, le deuxième étant propre, modulo bords de tels.
En tant qu'ensemble, il se définit aussi en termes de la vari té $X\mathrm{x}Y$
munie simplement de la famille des supports formée des ensembles fermés
propres sur Y ..... \handwritten{Les $B^{H}_{m}(X)$ (X compact disons) peuvent}
\handwritten{se dé\unsure{crire|finir} homotopiquement à la Thom, comme les groupes d'homotopie}
\handwritten{stables d'espaces $\underline{\mathrm{Hom}}(X, MH(N-m))$ - - - - -.}

\margin{\handwritten{48}}
